% Capítulo — TCC 1
\chapter{Estudo de caso e resultados preliminares}
\label{chap:estudodecaso}

Este capítulo descreve o contexto em que o sistema é avaliado: o Gov Hub, o
levantamento das junções entre fontes no código de produção e o repositório
do Ministério da Cultura, onde o sistema será aplicado. Termina com o estado
atual da implementação e os primeiros resultados.

\section{Gov Hub}
\label{sec:estudodecaso-govhub}

O Gov Hub é uma plataforma de dados de código aberto desenvolvida pelo Lab
Livre, laboratório de engenharia de software da UnB, com ministérios do
governo federal. Cada ministério tem um repositório
\texttt{data-application-*} na organização GovHub-br. A plataforma ingere
fontes federais como SIAFI, Transferegov, Compras.gov, SICONV, SALIC e IBGE
num PostgreSQL, com orquestração pelo Airflow, e mantém as junções entre elas
como \emph{models} dbt em produção.

\section{Inventário das junções entre fontes}
\label{sec:estudodecaso-inventario}

Para delimitar o problema, a autora leu o código dos repositórios do Gov Hub e
listou 29 condições de junção entre fontes diferentes no código de produção,
contando separadamente as chaves alternativas que uma junção tenta em
sequência. A cascata de dezesseis passos que extrai o número do TED de campos
de texto livre conta como uma condição. Uma delas está num \emph{model}
desativado e outra no código de ingestão. A Tabela~\ref{tab:inventario}
agrupa as condições por tipo de atrito. A classificação foi feita pela
autora, por leitura estática, e será auditada e congelada antes da avaliação.

\begin{table}[htbp]
\centering
\caption{Atritos nas 29 condições de junção entre fontes do código de produção
do Gov Hub (uma junção pode ter vários tipos; a junção sem transformação também
é contada em grafia)}
\label{tab:inventario}
\footnotesize
\begin{tabular}{@{}lrl@{}}
\toprule
\textbf{Atrito} & \textbf{Junções} & \textbf{Exemplo} \\
\midrule
Formato ou tipo & 15 & NE curta $\times$ identificador completo \\
Chave derivada & 9 & UG + gestão + NE $\rightarrow$ identificador SIAFI \\
Grafia ou nomenclatura & 5 & nomes de parlamentares, acentos \\
Extração de texto livre & 4 & número da transferência na descrição \\
Equivalência de código & 3 & a mesma unidade sob duas siglas \\
Estrutura aninhada & 1 & identificador dentro de JSON bruto \\
\midrule
Sem transformação & 1 & \\
\bottomrule
\end{tabular}
\end{table}

Apenas uma junção é uma igualdade direta entre colunas sem transformação. Ao
menos seis chaves são compostas, várias junções usam cascatas de
alternativas, e duas usam o CPF, um identificador pessoal, e ficam fora do
estudo. A única junção com estrutura aninhada é uma dessas duas, de modo que
esse atrito aparece apenas nos pares perturbados.

O inventário também mostra por que o catálogo e o Mapa de Fluxo são
necessários. A decomposição do identificador SIAFI (os seis primeiros
caracteres são a UG, os doze últimos são o documento) aparece em pelo menos
cinco \emph{models} e é a chave derivada mais frequente do domínio. As
junções que tiram a chave de texto livre (o número do TED na descrição da
nota de empenho, o número do convênio no texto da NE) só são reconhecíveis
por quem conhece o fluxo orçamentário.

\section{O repositório do Ministério da Cultura}
\label{sec:estudodecaso-minc}

O repositório \texttt{data-application-minc} mantém o \emph{pipeline} de
dados do MinC no Gov Hub. Ele ingere o Transferegov, o BB Gestão Ágil, o
SALIC, a Ancine, o SICONFI, o Banco Central, o Mapas Culturais e as bases do
IBGE (SIDRA, MUNIC e a base territorial FCU),
transforma os dados com o dbt em PostgreSQL e serve o acompanhamento das
metas da PNAB e da Lei Paulo Gustavo. O trabalho é financiado pelo TED n.º
01/2026/SGE/SE/MinC.

\subsection{Junções entre fontes no MinC}

Um levantamento estático no \emph{commit} \texttt{c588691} (10 de novembro de
2026), que percorre a linhagem de cada um dos \emph{models} dbt até as suas
fontes, encontrou apenas três \emph{models} que cruzam fontes diferentes:

\begin{itemize}
\item a junção das planilhas de contemplados do relatório de gestão com o
território do IBGE, por uma chave derivada de nome do município sem acento e
sigla da UF;
\item duas junções entre Ancine, BB Gestão Ágil, SALIC e as planilhas de
contemplados pelo documento do beneficiário, que pode ser CPF.
\end{itemize}

Só a primeira pode entrar no benchmark; as outras usam um identificador
pessoal. O restante dos cruzamentos do repositório acontece dentro de uma
mesma fonte, como o SALIC com o SALIC ou as edições da MUNIC entre si.
% TODO: auditar este levantamento (feito por script, em 15/11/2026) e
% registrar o script no repositório do sistema.

\subsection{Pares ainda não integrados}

O MinC tem, portanto, poucas junções em produção e várias fontes que
descrevem as mesmas entidades sem estar cruzadas. É nesse espaço que o
sistema pode gerar valor para o ministério: descobrir novos pares de
integração, testar a metodologia fora do benchmark e entregar os cruzamentos
confirmados como \emph{models}. São candidatos iniciais à QP4:

\begin{itemize}
\item Ancine $\times$ SALIC: a planilha da Ancine tem uma coluna com o número
SALIC do projeto, que corresponde ao PRONAC do SALIC;
\item SALIC $\times$ Transferegov e FCU: o SALIC usa o código IBGE do
município com seis dígitos, e as outras fontes, com sete; a ponte entre os
dois formatos existe numa tabela do próprio SALIC;
\item SICONFI e IBGE SIDRA $\times$ território da FCU, pelo código do
município.
\end{itemize}

A lista será revista com a equipe do Gov Hub no MinC antes do congelamento.
Esses pares não têm gabarito prévio, e por isso entram na QP4 e não no
benchmark da QP1.

\section{Resultados preliminares}
\label{sec:estudodecaso-preliminares}

\subsection{Cobertura do catálogo}

A autora conferiu, em 5 de novembro de 2026, cada uma das 29 condições contra
o rascunho do catálogo, por leitura estática do SQL. Catorze cabem em no
máximo três operações por lado; quatro precisam de \emph{extrair
identificador(tipo)} em pelo menos uma etapa da cascata; três ficam fora do
benchmark; e oito dependem de definições de operações ainda não congeladas.
Como o rascunho do catálogo saiu deste mesmo inventário, essa contagem
descreve o encaixe, não a generalização: os pares reais vão testar a escolha
da chave, não a cobertura do catálogo.

\subsection{Ilustração do protótipo}

Num par de desenvolvimento com bases do IBGE, o caminho só de regras do
protótipo (que não é o \emph{baseline} B2) escolheu \texttt{nome}
$\leftrightarrow$ \texttt{nome}. A chave correta é o código da UF, que é o
prefixo do código do município e também aparece aninhado num JSON: uma chave
derivada que a similaridade de nomes não encontra. O caso ilustra o problema;
não é evidência a favor do sistema proposto.

\subsection{Primeiro incremento por SDD}

Um incremento de delimitação do repositório, fora do protocolo pareado, foi
especificado com oito critérios de aceite e implementado por um agente de
código em duas iterações, sem intervenções. A especificação foi alterada duas
vezes depois de aprovada. A verificação encontrou uma divergência entre a
especificação e o código: a especificação supunha que o \emph{linter} já
passava, mas havia 78 erros. A divergência só apareceu quando os comandos de
verificação dos critérios foram executados, que é o papel que os
testes-oráculo escondidos assumem no protocolo pareado.

\subsection{Estado da implementação}

Existem a coleta de evidências e um caminho de decisão sem LLM. A chamada de
decisão tipada, a abstenção, o catálogo, o achatamento de JSON aninhado e a
renderização dbt ainda serão construídos, cada um como um incremento por SDD.
